\documentclass[11pt,a4paper]{article}

\usepackage[utf8]{inputenc}
\usepackage[T1]{fontenc}
\usepackage[ngerman]{babel}
\usepackage{amsmath,amssymb,mathtools}
\usepackage{booktabs}
\usepackage{geometry}
\usepackage{hyperref}
\usepackage{microtype}
\usepackage{siunitx}
\usepackage{xcolor}

\geometry{left=28mm,right=28mm,top=28mm,bottom=30mm}
\hypersetup{
  colorlinks=true,
  linkcolor=blue!55!black,
  urlcolor=blue!55!black,
  citecolor=blue!55!black,
  pdftitle={Gabriel--Hurwitz-Protokoll},
  pdfauthor={Thomas Hoffbauer}
}

\sisetup{
  locale = DE,
  output-decimal-marker = {,},
  group-separator = {.}
}

\newcommand{\Li}{\operatorname{Li}}
\newcommand{\piPrime}{\pi}
\newcommand{\HurwitzStates}{\mathcal{H}}
\newcommand{\Deviation}{\Delta}
\newcommand{\Koch}{B_{\mathrm{Koch}}}

\title{\textbf{Gabriel--Hurwitz-Protokoll}\\
\large Arithmetische Interferenz, Hurwitz-Quaternionen und der Eichpunkt \(x=1422\)}
\author{Dr. Thomas Hoffbauer}
\date{\today}

\begin{document}
\maketitle

\begin{abstract}
Dieses Protokoll verbindet das klassische Paradoxon von Gabriels Horn mit einer
gewichteten Primzahlgeometrie über Hurwitz-Quaternionen. Ausgangspunkt ist die
Fluktuation
\[
  \Deviation(x) = \piPrime(x)-\Li(x),
\]
also die Abweichung der realen Primzahlzählfunktion vom glatten
Integrallogarithmus. Numerisch zeigt sich im Bereich \(1000 \leq x < 3000\) ein
kritischer Eichpunkt bei \(x=1422\). Dort liegt die lokale Auslastung der
verwendeten Rigiditätsskala bei \(\num{119.92}\,\%\). Im Modell wird diese
Kompression durch eine Primzahllücke zwischen \(1409\) und \(1423\) und durch
die gewichtete Hurwitz-Masse \(24p\) der angrenzenden Primzahlschalen
interpretiert.
\end{abstract}

\section{Gabriels Horn als Kontinuumsmodell}

Gabriels Horn entsteht, wenn der Graph
\[
  y=\frac{1}{x}, \qquad x \geq 1,
\]
um die \(x\)-Achse rotiert. Das Volumen dieses unendlich langen Körpers ist
endlich:
\[
  V=\pi\int_{1}^{\infty}\frac{1}{x^2}\,\mathrm{d}x=\pi.
\]
Die Mantelfläche divergiert dagegen:
\[
  A=2\pi\int_{1}^{\infty}\frac{1}{x}\sqrt{1+\frac{1}{x^4}}\,\mathrm{d}x
  =\infty.
\]
Damit liefert Gabriels Horn eine präzise mathematische Metapher: Ein
kontinuierliches System kann global endlich und lokal dennoch unerschöpflich
sein. Genau diese Spannung zwischen Kontinuum und diskreter Struktur tritt in
der Primzahlgeometrie erneut auf.

\section{Primzahltrend und arithmetische Fluktuation}

Die glatte Näherung an die Primzahlzählfunktion wird durch den
Integrallogarithmus beschrieben:
\[
  \Li(x)=\int_{2}^{x}\frac{1}{\log t}\,\mathrm{d}t.
\]
Die reale Primzahlzählung \(\piPrime(x)\) ist hingegen eine Treppenfunktion.
Die zentrale Messgröße des Modells ist deshalb
\[
  \Deviation(x)=\piPrime(x)-\Li(x).
\]
Als Vergleichsskala verwenden wir die im Skript getestete von-Koch-artige
Rigiditätsskala
\[
  \Koch(x)=\frac{1}{8\pi}\sqrt{x}\log x.
\]
Die normierte Kompression ist
\[
  R(x)=\frac{|\Deviation(x)|}{\Koch(x)}.
\]
Sie misst nicht eine bewiesene globale Schranke, sondern die lokale Auslastung
der im Modell verwendeten Rigiditätsskala.

\section{Numerischer Eichpunkt}

Im Suchfenster \(1000\leq x<3000\) ergibt die Diagnose den maximalen
Kompressionspunkt
\[
  x_\ast=1422.
\]
Die zugehörigen Kennzahlen lauten:

\begin{table}[htbp]
\centering
\caption{Kritischer Eichpunkt des Bamberger Modells}
\begin{tabular}{@{}lr@{}}
\toprule
Größe & Wert \\
\midrule
Suchfenster & \(1000 \leq x < 3000\) \\
Eichpunkt \(x_\ast\) & \(1422\) \\
Abweichung \(|\Deviation(x_\ast)|\) & \(\num{13.0627}\) \\
Rigiditätsskala \(\Koch(x_\ast)\) & \(\num{10.8927}\) \\
Auslastung \(R(x_\ast)\) & \(\num{119.92}\,\%\) \\
\bottomrule
\end{tabular}
\end{table}

Der Wert \(x=1422\) liegt unmittelbar vor der nächsten Primzahl \(1423\).
Dadurch bleibt \(\piPrime(x)\) lokal konstant, während \(\Li(x)\) stetig
weiterwächst. Die Kompression erreicht daher genau am Ende der primzahlfreien
Zone ihr lokales Maximum.

\section[Kritisches Fenster 1400 bis 1450]{Kritisches Fenster \(1400\leq x\leq 1450\)}

Die Primzahlen im kritischen Fenster sind
\[
  1409,\ 1423,\ 1427,\ 1429,\ 1433,\ 1439,\ 1447.
\]
Die lokalen Lücken lauten:

\begin{table}[htbp]
\centering
\caption{Lokale Lückenstruktur im kritischen Fenster}
\begin{tabular}{@{}ccc@{}}
\toprule
Vorherige Primzahl & Nächste Primzahl & Lücke \\
\midrule
1409 & 1423 & 14 \\
1423 & 1427 & 4 \\
1427 & 1429 & 2 \\
1429 & 1433 & 4 \\
1433 & 1439 & 6 \\
1439 & 1447 & 8 \\
\bottomrule
\end{tabular}
\end{table}

Die maximale arithmetische Spannung im Fenster ist somit die Lücke
\[
  1423-1409=14.
\]
Im Bild des Modells ist das Intervall \(1410,\dots,1422\) ein lokales
arithmetisches Vakuum: kein neuer Primzahlknoten tritt auf, obwohl das
kontinuierliche Hintergrundpotential weiter anwächst.

\section{Hurwitz-Gewichtung}

Für eine Primzahl \(p\) wird die algebraische Zustandszahl der zugehörigen
Hurwitz-Sphäre im Modell durch
\[
  \HurwitzStates(p)=24p
\]
gewichtet. Damit wird aus der bloßen Zählung von Primzahlen eine gewichtete
Massendichte des Gitters.

Für die globale Schnittstelle bis \(x=2000\) ergibt die Diagnose:

\begin{table}[htbp]
\centering
\caption{Quaternionische Feld-Schnittstelle bis \(x=2000\)}
\begin{tabular}{@{}lr@{}}
\toprule
Größe & Wert \\
\midrule
Anzahl der Primzahlknoten & \(303\) \\
Kontinuierliches Hintergrundpotential \(\Li(2000)\) & \(\num{313.7641}\) \\
Gesamtzahl der Hurwitz-Zustände & \(\num{6649200}\) \\
Globale Energiedichte & \(\num{21191.72}\) \\
\bottomrule
\end{tabular}
\end{table}

Die lokalen Massenpunkte um den Eichpunkt sind:

\begin{table}[htbp]
\centering
\caption{Lokale Hurwitz-Massen im kritischen Fenster}
\begin{tabular}{@{}ccc@{}}
\toprule
Primzahl \(p\) & Hurwitz-Zustände \(24p\) & Anteil an der Gesamtmasse \\
\midrule
1409 & 33816 & \(\num{0.5086}\,\%\) \\
1423 & 34152 & \(\num{0.5136}\,\%\) \\
1427 & 34248 & \(\num{0.5151}\,\%\) \\
1429 & 34296 & \(\num{0.5158}\,\%\) \\
\bottomrule
\end{tabular}
\end{table}

Die 14er-Lücke zwischen \(1409\) und \(1423\) ist damit nicht nur eine Lücke in
der Primzahlzählung, sondern ein massenfreies Intervall zwischen zwei stark
gewichteten Hurwitz-Knoten.

\section{Fernfeldkontrolle}

Als Fernfeld-Stichprobe wurde \(x=100000\) berechnet:

\begin{table}[htbp]
\centering
\caption{Fernfeld-Stichprobe}
\begin{tabular}{@{}lr@{}}
\toprule
Größe & Wert \\
\midrule
\(\piPrime(100000)\) & \(9592\) \\
\(\Li(100000)\) & \(\num{9628.76}\) \\
\(\Deviation(100000)\) & \(\num{-36.76}\) \\
\(\Koch(100000)\) & \(\pm\num{144.86}\) \\
Auslastung & \(\num{25.38}\,\%\) \\
\bottomrule
\end{tabular}
\end{table}

Im Fernfeld ist die relative Auslastung der Rigiditätsskala deutlich kleiner.
Das unterstützt die Modellinterpretation, dass der Bereich um \(x=1422\) ein
lokaler Eichpunkt des Nahfeldes ist.

\section{Resonanzdiagnose}

Zur groben Orientierung wurde außerdem eine abgeschnittene Zeta-Welle aus den
ersten drei nichttrivialen Nullstellen-Frequenzen verwendet:
\[
  W(x)=\sum_{\gamma\in\{14.134725,\,21.022040,\,25.010858\}}
  \frac{2\sin(\gamma\log x)}{\gamma\sqrt{x}}.
\]
Die Werte lauten:
\[
  W(1500)=\num{0.0029},\qquad W(8000)=\num{0.0011}.
\]
Diese Größe ist nicht die Fluktuation selbst, sondern eine normierte
Resonanzsonde für die im Plot sichtbaren Oszillationsanteile.

\section{Modellfazit}

Das Gabriel--Hurwitz-Protokoll lässt sich in drei Schritten lesen:

\begin{enumerate}
  \item Gabriels Horn liefert das Kontinuumsbild: endliches Volumen bei
  unendlicher Oberfläche.
  \item Der Integrallogarithmus \(\Li(x)\) liefert den glatten Primzahlstrom,
  während \(\piPrime(x)\) die diskrete Treppenstruktur zeigt.
  \item Die Hurwitz-Gewichtung \(24p\) macht aus Primzahlen gewichtete
  Gitterknoten. Die Lücke \(1409\to1423\) erzeugt dadurch im Modell ein lokales
  massenfreies Intervall zwischen zwei schweren algebraischen Zustandsknoten.
\end{enumerate}

Der Eichpunkt \(x=1422\) ist damit der numerisch isolierte Ort maximaler lokaler
Kompression im untersuchten Nahfeld. Er erklärt, warum die Kurve
\(\piPrime(x)-\Li(x)\) dort besonders stark an die gewählte Rigiditätsskala
heranläuft und anschließend wieder in den stabileren Korridor zurückkehrt.

\end{document}